\cotefile{62}
\scanpage{68}
Expression au moyen de la fonction $\sigma$
\[
\boxed{\mathrm{sn}\,z=\sigma\omega_1\,e^{\eta_1 z}\frac{\sigma z}{\sigma(z+\omega_1)}
=-\sigma\omega_1\,e^{-\eta_1 z}\frac{\sigma z}{\sigma(z-\omega_1)}} \tag*{(6)}
\]

\unsure{On a|On|$q$} \deleted{$\mathrm{sn}\,z$ admet $\sigma(z$}

On en déduit
\[
\left\{\begin{aligned}
\mathrm{sn}(z+2\omega_1)&=\mathrm{sn}\,z\\
\mathrm{s}\deleted{\unsure{m|n}}\mathrm{n}(z+2\omega_2)&=-\mathrm{sn}\,z\\
\mathrm{sn}(z+2\omega_3)&=-\mathrm{sn}\,z\\
\text{d'où}\qquad\mathrm{sn}(z+4\omega_2)&=\mathrm{sn}\,z\\
\mathrm{sn}(z+4\omega_3)&=\mathrm{sn}\,z
\end{aligned}\right. \tag*{(7)}
\]
Donc $\mathrm{sn}\,z$ admet les périodes $2\omega_1$ et $4\omega_2$ (et $4\omega_3$)

Ses poles sont $z=\omega_1+2n_1\omega_1+4n_{\unsure{1|2}}\omega_2$, ses zéros $z=2n_1\omega_1+4n_1\omega_2$, poles et zéros sont simples ; $\mathrm{sn}\,z$ est une fonction d'ordre 2.

Ses zéros sont homologues de $0$ ou $2\omega_{2}$, où la dérivée est respectivement $1$ et $-1$ ; les poles sont homologues de pole $\omega_1$ et $\omega_1+2\omega_2$, dont les résidus sont resp\textsuperscript{t} $\frac{1}{k}$ et $-\frac{1}{k}$ (en supposant toujours choisi \inserted{le signe de} $k$ tel que l'on ait (5), ou, ce qui revient au même $\displaystyle k=-\frac{e^{\eta_1\omega_1}}{\sigma^{\unsure{2|1|\prime}}\omega_1}$)

On voit sur 3 ou 4 que $\mathrm{sn}\,z$ est paire ou impaire ; elle n'est pas \deleted{im}paire, car on aurait $\mathrm{sn}\,\omega_2=\mathrm{s}\deleted{\unsure{n}}\mathrm{n}\,{-\omega_2}$, mais aussi $\mathrm{sn}\,\omega_2=-\mathrm{sn}\,{-\omega_2}$, d'où $\mathrm{sn}\,\omega_2=0$, ce qui n'est pas. On peut aussi s'en assurer sur 6
\[
\boxed{\mathrm{sn}\,{-z}=-\mathrm{sn}\,z}
\]

Sur 3 ou 4, on voit $\mathrm{sn}^2\omega_2=\frac{1}{\deleted{\unsure{k^2|k}}}$ ; en choisissant convenablement le signe de $\omega_2$, on aura \uline{$\mathrm{sn}\,\omega_2=\frac{1}{\deleted{\unsure{k}}}$} ; l'équation différentielle vérifiée par $\mathrm{sn}\,z$ donne \uline{$\mathrm{sn}'\omega_2=0$} de $=$ $\mathrm{sn}\,{-\omega_2}=-1\quad\mathrm{sn}'{-\omega_2}=0$.

On voit alors sur 4 \deleted{que}, puis sur l'équation différentielle, que $\mathrm{sn}(\omega_1+\omega_2)=\frac1k$ soit \uline{$\mathrm{sn}\,\omega_3=-\frac1k\quad\mathrm{sn}'\omega_3=0$}.

% archive sheet number: 67
\scanpage{69}
On pose
\[
\left\{\begin{aligned}
\mathrm{cn}^2z&=1-\mathrm{sn}^2z=\frac{\unsure{\wp|p}z-e_2}{\wp z-e_1}=\deleted{\unsure{-\wp(z-\omega_2)|-\wp(z-\omega_1)}}-\frac1{k^2}\bigl[\wp(z-\omega_1)-e_3\bigr]
&\boxed{\mathrm{cn}\,0=+1}\\
\mathrm{dn}^2z&=1-k^2\mathrm{sn}^2z=\frac{\wp z-e_3}{\wp z-e_1}=-\bigl[\wp(z-\omega_1)-e_2\bigr]
&\boxed{\mathrm{dn}\,0=+1}
\end{aligned}\right. \tag*{8}
\]

On a
\[
\left\{\begin{aligned}
\mathrm{cn}\,z&=\frac{\sigma\omega_1}{\sigma\omega_2}e^{(\eta_{2}-\eta_1)z}\frac{\sigma(z-\omega_2)}{\sigma(z-\omega_1)}
=\frac{\sigma\omega_1}{\sigma\omega_2}e^{-(\eta_2-\eta_1)z}\frac{\sigma(z+\omega_2)}{\sigma(z+\omega_1)}\\
\mathrm{dn}\,z&=\frac{\sigma\omega_1}{\sigma\omega_3}e^{(\eta_3-\eta_1)z}\frac{\sigma(z-\omega_3)}{\sigma(z-\omega_1)}
=\frac{\sigma\omega_1}{\sigma\omega_3}e^{-(\eta_3-\eta_1)z}\frac{\sigma(z+\omega_3)}{\sigma(z+\omega_1)}
\end{aligned}\right. \tag*{9}
\]

\[
\left\{\begin{aligned}
&{\times}\ \mathrm{cn}\,{-z}=\mathrm{cn}\,z &&\mathrm{dn}\,{-z}=\mathrm{dn}\,z &&\mathrm{cn}\,z\text{ et }\mathrm{dn}\,z\text{ sont paires}\\
&{\times}\left\{\begin{aligned}
\mathrm{cn}(z+2\omega_1)&=-\mathrm{cn}\,z\\
\mathrm{cn}(z+2\omega_2)&=-\mathrm{cn}\,z
\end{aligned}\right. &&
\begin{aligned}\mathrm{dn}(z+2\omega_3)&=-\mathrm{dn}\,z\\
\mathrm{dn}(z+2\omega_1)&=-\mathrm{dn}\,z\end{aligned}\\
&{\times}\left\{\begin{aligned}
\mathrm{cn}(z+2\omega_3)&=\mathrm{cn}\,z\\
\mathrm{cn}(z+4\omega_1)&=\mathrm{cn}\,z\\
\mathrm{cn}(z+4\omega_2)&=\mathrm{cn}\,z
\end{aligned}\right. &&
\begin{aligned}\mathrm{dn}(z+2\omega_2)&=\mathrm{dn}\,z\\
\mathrm{dn}(z+4\omega_3)&=\mathrm{dn}\,z\\
\mathrm{dn}(z+4\omega_1)&=\mathrm{dn}\,z\end{aligned} &&
\begin{aligned}&\mathrm{cn}\text{ admet pour périodes}\\&2\omega_3,4\omega_1,4\omega_2\\
&\mathrm{dn}\text{ admet pour période}\\&2\omega_3,4\omega_{\unsure{1|3}},4\omega_2\end{aligned}
\end{aligned}\right. \tag*{10}
\]

$\mathrm{cn}\,z$ et $\mathrm{dn}\,z$ sont fonctions elliptiques d'ordre 2 ;

\uuline{Poles de $\mathrm{cn}\,z$} homologues de $\omega_1$, résidu $\frac{i}{k}$ \inserted{si $\omega_1,\omega_2$ \deleted{est direct}}, et \deleted{$\omega_{\unsure{1|2}}$} \deleted{\unsure{(orienté) vers $2\omega_2$|(si $\omega_2$) vers $2\omega_2$}} \inserted{$-\omega_1$} résidu $-\frac{i}{k}$

(l'ambiguité de signe est levée par 9, joint à 6)

\deleted{Poles} \uuline{Zéros de $\mathrm{cn}\,z$} homologues de $+\omega_2$ et $-\omega_2$ ; on a,
\[
\mathrm{cn}'\omega_2=-\frac{e^{\eta_2\omega_2}}{\sigma^2\omega_2}\qquad (9,\text{ puis }6)
\]
; a partir de relation de définition de $\mathrm{cn}\,z$, on montre $\mathrm{cn}^{\prime\,2}\omega_2=-\mathrm{sn}''\omega_2$
\[
\mathrm{cn}'\omega_2=\pm\sqrt{1-k^2}
\]
Enfin
\[
\mathrm{cn}\,\omega_3=\frac{i}{k}\frac{e^{\eta_2\omega_2}}{\sigma^2\omega_2}=-\frac{i}{k}\mathrm{cn}'\omega_2=\pm\frac{i}{k}\sqrt{1-k^2}\qquad \mathrm{cn}'\omega_3=0
\]

% archive sheet number: 68
\scanpage{70}
\uuline{Poles de $\mathrm{dn}\,z$} homologues de \uline{$\omega_1$, résidu $i$}, et $-\omega_1$, résidu $-i$

\uuline{\deleted{\unsure{Poles|Zéros}}Zéros de $\mathrm{dn}\,z$} homologues de \uline{$\omega_3$, dérivée $k\,\mathrm{cn}\,\omega_3$}$=i\,\frac{e^{\eta_2\omega_2}}{\sigma^2\omega_2}=-i\,\mathrm{cn}'\omega_2=\pm i\sqrt{1-k^2}$ et de $-\omega_3$, (dérivées symétriques)

Enfin $\mathrm{dn}\,\omega_2=-\mathrm{cn}'\omega_2=\frac{e^{\eta_2\omega_2}}{\sigma^2\omega_2}=\pm\sqrt{1-k^2}\qquad\mathrm{dn}'\omega_2=0$

\uline{formules de dérivation}
\[
\left\{\begin{aligned}
\mathrm{sn}'z&=\mathrm{cn}\,z\,\mathrm{dn}\,z\\
\mathrm{cn}'z&=-\mathrm{sn}\,z\,\mathrm{dn}\,z\\
\mathrm{dn}'z&=-k^2\mathrm{sn}\,z\,\deleted{\unsure{d}}\mathrm{cn}\,z
\end{aligned}\right. \tag*{11}
\]

formules liant $\mathrm{sn}\,z$, $\deleted{\unsure{d}}\mathrm{cn}\,z$, $\mathrm{dn}\,z$ deux à deux
\[
\boxed{\mathrm{cn}^2z+\mathrm{sn}^2z=1\qquad
\deleted{\unsure{k^2}}\mathrm{dn}^2z+k^2\mathrm{sn}^2z=1\qquad
\mathrm{dn}^2z-k^2\mathrm{cn}^2z=1-k^2} \tag*{(12)}
\]

$\mathrm{cn}\,z$ et $\mathrm{dn}\,z$ vérifient des équations analogues à celle de $\mathrm{sn}\,z$
\[
\left\{\begin{aligned}
\mathrm{cn}^{\prime\,2}z&=(1-\mathrm{cn}^2z)(1-k^2\unsure{-|+}k^2\mathrm{cn}^2z)=(1-k^2)(1-\mathrm{cn}^2z)\left(1+\frac{k^2}{1-k^2}\mathrm{cn}^2z\right)\\
\mathrm{dn}^{\prime\,2}z&=(1-\mathrm{dn}^2z)(\mathrm{dn}^2z-(1-k^2)=-(1-k^2)(1-\mathrm{dn}^2z)\left(1-\frac{1}{1-k^2}\mathrm{dn}^2z\right)
\end{aligned}\right. \tag*{(13)}
\]

Aussi $\mathrm{cn}\,z$ et $\mathrm{dn}\,z$ s'expriment\deleted{\unsure{is}}\unsure{-il|-ils} directement au moyen de la fonction $\mathrm{sn}\,z$ qui correspond respectivement à la solution de l'équation $\left(\frac{dZ}{dz}\right)^2=(1-Z^2)(1-{\unsure{\rho|p}}^{2}Z^2)$ \deleted{\unsure{avec|ou|autre}} \unsure{pour|avec} $\rho=\frac{k}{\sqrt{1-k^2}}$ et $\rho=\frac{1}{\sqrt{1-k^2}}$

% archive sheet number: 69
\scanpage{71}
\[
\mathrm{sn}_1^2z=\frac1{\unsure{\wp|p}z-e_1}\qquad(\mathrm{sn}'_1\,0=+1) \tag*{(1)}
\]
\[
\mathrm{sn}_1z=\sigma\omega_1\,e^{\eta_1z}\frac{\sigma z}{\sigma(z+\omega_1)}
=-\sigma\omega_1\,e^{-\eta_1z}\frac{\sigma z}{\sigma(z-\omega_1)} \tag*{(2)}
\]
\[
\deleted{\mathrm{sn}_1^2z=\frac{1}{\wp z-e_1}}
\]
\[
\begin{aligned}
\mathrm{sn}_1z\,\mathrm{sn}_1(z+\omega_1)&=\frac1{k_1}&
\mathrm{sn}_1^2z&=\frac1{k_1^2}\bigl(\wp(z-\omega_1)-e_1\bigr)&
k_1&=-\frac{e^{\eta_1\omega_{\unsure{1|2}}}}{\sigma^2\omega_1}\\
&&&&k_1^2&=(e_3-e_1)(e_{\unsure{3|2}}-e_1)
\end{aligned}\tag*{(3)}
\]
\[
\left\{\begin{aligned}
\mathrm{sn}_1(z+2\omega_1)&=\mathrm{sn}_1z\\
\mathrm{sn}_1(z+2\omega_2)&=-\mathrm{sn}_1z\\
\mathrm{sn}_1(z+2\omega_3)&=-\mathrm{sn}_1z\\
\mathrm{sn}_1(z+4\omega_{\unsure{2|1}})&=\mathrm{sn}_1z\\
\mathrm{sn}_1(z+4\omega_3)&=\mathrm{sn}_1z\\
\mathrm{sn}_1\,{-z}&=-\mathrm{sn}_1z\\
\mathrm{sn}_1(2\omega_2-z)&=\mathrm{sn}_1z\\
\mathrm{sn}_1(2\omega_3-z)&=\mathrm{sn}_1z
\end{aligned}\right. \tag*{(4)}
\]
$\mathrm{sn}_1z$ a les périodes $2\omega_1,4\omega_2,4\omega_3$.

La substitution $z\to z+2\omega_2$ et $z\to z+2\omega_3$ équivaut à un changement de signe

$\mathrm{sn}_1z$ est impaire

Les substitutions $z\to2\omega_2-z$ et $z\to2\omega_3-z$ laissent $\mathrm{sn}_1z$ invariante

$\mathrm{sn}_1z$, fonction elliptique ayant les périodes $2\omega_1,4\omega_2,4\omega_3$, d'ordre deux

\uline{poles de $\mathrm{sn}_1z$} les homologues de \uline{$\omega_1$, résidu $+\frac1{k_1}$}, et $\omega_1+2\omega_2$ ou $\omega_1+2\omega_3$, résidu $-\frac1{k_1}$

\uline{zéros de $\mathrm{sn}_1z$} homologues de \uline{$0$, dérivée $+1$}, et $2\omega_2$ ou $2\omega_3$, dér. $-1$
\[
\left\Vert\begin{aligned}
\mathrm{sn}_1\omega_2&=\frac{\sigma\omega_1\,\sigma\omega_2}{\sigma(\omega_1+\omega_2)}e^{\eta_1\omega_2}=\pm\frac{1}{\sqrt{e_2-e_1}}&
\mathrm{sn}'_1\omega_2&=0\\
\mathrm{sn}_1\omega_3&=\frac{\sigma\omega_1\,\sigma\omega_3}{\sigma(\omega_1+\omega_{\unsure{2|3}})}e^{\eta_1\omega_3}=\pm\frac{1}{\sqrt{e_{\unsure{3|2}}-e_1}}&
\mathrm{sn}'_1\omega_3&=0
\end{aligned}\right.
\]

% archive sheet number: 70
\scanpage{72}
On pose
\[
\left\{\begin{aligned}
\mathrm{cn}_1^2z&=1-(e_2-e_1)\mathrm{sn}_1^2z=\frac{\unsure{\wp|p}z-e_2}{\wp z-e_1}&(\mathrm{cn}_1\,0=1)\\
\mathrm{dn}_1^2z&=1-(e_3-e_1)\mathrm{sn}_1^2z=\frac{\wp z-e_3}{\wp z-e_1}&(\mathrm{dn}_1\,0=1)
\end{aligned}\right.\tag*{5}
\]
\[
\left\{\begin{aligned}
\mathrm{cn}_1z&=\frac{\sigma\omega_1}{\sigma\omega_2}e^{(\eta_2-\eta_1)z}\frac{\sigma(z-\omega_2)}{\sigma(z-\omega_1)}
=\frac{\sigma\omega_1}{\sigma\omega_2}e^{-(\eta_2-\eta_1)z}\frac{\sigma(z+\omega_2)}{\sigma(z+\omega_1)}\\
\mathrm{dn}_1z&=\frac{\sigma\omega_1}{\sigma\omega_3}e^{(\eta_3-\eta_1)z}\frac{\sigma(z-\omega_3)}{\sigma(z-\omega_1)}
=\frac{\sigma\omega_1}{\sigma\omega_3}e^{-(\eta_3-\eta_1)z}\frac{\sigma(z+\omega_3)}{\sigma(z+\omega_1)}
\end{aligned}\right.\tag*{6}
\]

\deleted{\uline{Poles de $\mathrm{cn}_1z$}}
\[
\left\{\begin{aligned}
\mathrm{cn}_1\,{-z}&=\mathrm{cn}\,z &\mathrm{dn}_1\,{-z}&=\mathrm{dn}_1z\\
\mathrm{cn}_1(z+2\omega_1)&=-\mathrm{cn}\,z &\mathrm{dn}_1(z+2\omega_1)&=-\mathrm{dn}_1z\\
\mathrm{cn}_1(z+2\omega_2)&=-\mathrm{cn}\,z &\mathrm{dn}_1(z+2\omega_{3})&=-\mathrm{dn}_1z\\
\Vert\ \mathrm{cn}_1(z+2\omega_3)&=\mathrm{cn}\,z &\mathrm{dn}_1(z+2\omega_2)&=\mathrm{dn}_1z\\
\mathrm{cn}_1(z+\deleted{2}4\omega_1)&=\mathrm{cn}\,z &\mathrm{dn}_1(z+4\omega_3)&=\mathrm{dn}_1z\\
\mathrm{cn}_1(z+4\omega_2)&=\mathrm{cn}\,z &\mathrm{dn}_1(z-4\omega_1)&=\mathrm{dn}_1z
\end{aligned}\right.\tag*{7}
\]
$\mathrm{cn}_1z$ et $\mathrm{dn}_1z$ sont paires

$\mathrm{cn}_1z$ et $\mathrm{dn}_1z$ ont resp\textsuperscript{t} les périodes $2\omega_3,4\omega_1,4\omega_3$ et $2\omega_2,4\omega_3,4\omega_{\unsure{1|2}}$ ; elles sont d'ordre deux.

\uline{Poles de $\mathrm{cn}_1z$} : homologues de $\omega_1$, \deleted{résidu $\unsure{z\omega_1+2\omega_3|\omega_1+2\omega_3|3\omega_1+2\omega_3}$} \inserted{et} $-\omega_1$

Si la disposition des vecteurs $\omega_1,\omega_2$ est directe, le résidu en $\omega_1$ est
\[
\text{résidu en }\omega_1:\ \frac{i}{k_1\mathrm{sn}_1\omega_2}\ ;=\pm i\sqrt{e_3-e_1}
\]
\uline{Zéros de $\mathrm{cn}_1z$} homologues de $+\omega_2$ et $-\omega_2\deleted{+\unsure{2\omega_1}}$ \deleted{\unsure{ou $\omega_2+2\omega_3$|ou $\omega_1+2\omega_3$}} ;

\uline{Poles de $\mathrm{dn}_1z$} : homologues de $\omega_1$ et $-\deleted{\unsure{3|2}}\omega_1$ ; avec l'hypothèse précédente, on a
\[
\text{résidu en }\omega_1:\ \frac{-i}{k_1\mathrm{sn}_1\omega_3}=\pm i\sqrt{e_2-e_1}
\]
\uline{Zéros de $\mathrm{dn}_1z$} : homologues de $\omega_3$ et $-\omega_3\deleted{+\unsure{2\omega_2}}$

% archive sheet number: 71
\scanpage{73}
\margin{\unsure{$d$|$J$|$\delta$}}

\[
\boxed{\mathrm{sn}^{\prime\,2}z=(1-z^2)(1-k^2z^2)}\qquad \mathrm{sn}\,0=0\quad \mathrm{sn}'0=1 \tag*{(1)}
\]
\[
\boxed{\mathrm{sn}^2z=\frac{1}{\unsure{\wp|p}z-e_1}=\frac1{k^2}\bigl[\wp(z-\omega_1)-e_1\bigr]}\tag*{(2)}
\]
\[
\left\{\begin{aligned}
e_1&=-\frac{1+k^2}{3}\\
e_2&=\frac{2-k^2}{3}\\
e_3&=\frac{2k^2-1}{3}
\end{aligned}\right.
\qquad\begin{aligned}e_2-e_1&=1\\e_3-e_1&=k^2\end{aligned}
\]
\[
\boxed{\mathrm{sn}\,z=\sigma\omega_1\,e^{\eta_1 z}\frac{\sigma z}{\sigma(z+\omega_1)}
=-\sigma\omega_1\,e^{-\eta_1z}\frac{\sigma z}{\sigma(z-\omega_1)}}\tag*{(\unsure{1|3})}
\]
Si on choisit \uline{$k=-\frac{e^{\eta_1\omega_1}}{\sigma^2\omega_1}$} alors
\[
\boxed{\mathrm{sn}\,z\,\mathrm{sn}(z+\omega_1)=\frac1k}\tag*{(4)}
\]
\[
\left\{\begin{aligned}
\mathrm{sn}(z+2\omega_1)&=\mathrm{sn}\,z\\
\mathrm{sn}(z+2\omega_2)&=\mathrm{sn}\,z\\
\mathrm{sn}(z+2\omega_3)&=-\mathrm{sn}\,z\\
\mathrm{sn}(z+\deleted{2}4\omega_2)&=\mathrm{sn}\,z\\
\mathrm{sn}(z+4\omega_3)&=\mathrm{sn}\,z\\
\mathrm{sn}\,{-z}&=-\mathrm{sn}\,z
\end{aligned}\right.\tag*{(5)}
\]
$2\omega_1,4\omega_2$ et $4\omega_3$ sont périodes

$\mathrm{sn}\,z$ est impaire

La substitution $z\to z+2\omega_2$ ou $z+2\omega_{\unsure{3|2}}$ donne un changement de signe

\[
\left\{\begin{aligned}
&\text{\uuline{poles} : les homologues de }\text{\uuline{$\omega_1$}},\ \text{résidu }\frac1k\ ;\ \text{et }\omega_1+2\omega_2,\ \text{résidu }-\frac1k\\
&\text{\uuline{zéros} : les homologues de }\text{\uuline{$0$}},\ \text{dérivée }+1\ ;\ \text{et }2\omega_2,\ \text{dérivée }-1
\end{aligned}\right.\tag*{6}
\]
en choisissant convenablement $\omega_{2}$ :
\[
\underline{\mathrm{sn}\,\omega_2=1\quad\mathrm{sn}'\omega_2=0}
\qquad\underline{\mathrm{sn}\,\omega_3=\frac{-1}{k}\quad\mathrm{sn}'\omega_3=0}
\]

% archive sheet number: 72
\scanpage{74}
\[
\left\{\begin{aligned}
\mathrm{sn}'_1z&=\mathrm{cn}_1z\,\mathrm{dn}_1z\\
\mathrm{cn}'_1z&=-(e_{\unsure{2|1}}-e_1)\mathrm{sn}_1z\,\mathrm{dn}_1z\\
\mathrm{dn}'_1z&=-(e_3-e_1)\mathrm{sn}_1z\,\mathrm{cn}_1z
\end{aligned}\right.\tag*{8}
\]
$\mathrm{sn}'_1z$, $\mathrm{cn}'_1z$, $\mathrm{dn}'_1z$ vérifient des équations différentielles de type analogue

On a notamment
\[
\mathrm{sn}_1^{\prime 2}z=\bigl(1-(e_2-e_1)\mathrm{sn}_1^2z\bigr)
\bigl(1-(e_3-e_1)\mathrm{sn}_1^2z\bigr)\tag*{(9)}
\]

% archive sheet number: 73
\scanpage{75}
\begin{center}\fbox{Fonctions elliptiques}\end{center}

\inserted{\uline{Expression \unsure{des|en} fonctions elliptiques par $\unsure{\wp|p}u,\wp' u$}}

1) \uuline{Généralités}. On appelle \uline{fonction elliptique} une fonction méromorphe $f$ dans $C$ qui admet deux périodes $\omega_1,\omega_2$ linéairement ind. sur $R$ ; si $f$ n'est pas constante, son groupe des périodes, (sous-groupe fermé de $C=R^2$) est discret, soit $\unsure{\mathfrak P|\mathfrak B}$\inserted{ et de rang 2.} On étudie \inserted{l'ensemble} des fonctions elliptiques qui admettent $\mathfrak P$ comme groupe des périodes.

\deleted{\unsure{Soit donc|Soit la}} Soit $\{2\omega_1,2\omega_2\}$ une base de $\mathfrak P$ sur $Z$ (définie à une substitution entière de déterminant $\pm1$ près) on étudie l'ens des fonctions méromorphes telles que
\[
f(z+2\omega_1)=f(z+2\omega_2)=f(z)\tag*{(1)}
\]
Il forme un \uline{corps} $E$, qu'on peut évidemment identifier au corps des fonctions méromorphes sur la variété holomorphe quotient $C/\mathfrak P$, -- qui est topologiquement isomorphe au tore. -- On sait, d'après \deleted{\unsure{le|la}} \unsure{théorie|théorème} \deleted{\unsure{finitiste|finitude|générale}} \inserted{de Riemann}, que \deleted{\unsure{pour}} \inserted{pour} \deleted{toute} $f\in E$ non constante \deleted{\unsure{on}} (et il en existe) \deleted{\illegible{1}} $E$ est de degré fini sur $C(f)$, i. e \deleted{\unsure{toute}} fonction elliptique \inserted{$g\in E$} est \deleted{\unsure{égale}} solution d'une équation $P(f,g)=0$, $P$ polynôme en $f$ et $g$ \inserted{etc.} On peut ici obtenir élémentairement des résultats plus précis.

a') \deleted{\unsure{Existence d'une fon}} \uline{Une fonction elliptique sans pôle est une constante}

-- \uline{immédiat par} Liouville

\margin{\unsure{a)|a')} Somme des résidus nulle ; nb de zéros \unsure{égale|égal au} nb \unsure{des|de} pôles. \unsure{corollaire} : au moins deux \unsure{distincts|d'ordre}.}

b) Existence d'une fonction elliptique \unsure{non constante}
\[
\wp(z)=\frac{1}{z^2}+\sum_{m,n}'\left(\frac{1}{(z-\unsure{m|2m}\omega_1\unsure{+|-}n\omega_2)^2}-\frac{1}{(m\omega_1+n\omega_2)^2}\right)\tag*{(2)}
\]
définit une fonction méromorphe dans \deleted{\unsure{tout}} le plan, ayant \deleted{pour pôles} $\mathfrak P$ comme ensemble des pôles (pôles doubles)

On a
\[
\wp'(z)=-\frac{2}{z^3}-2\sum_{m,n}'\frac{1}{(z\unsure{+|-}m\omega_1+n\omega_2)^3}
=-2\sum\frac{1}{(z-2m\omega_1\unsure{-|+}2n\omega_2)^3}\tag*{3)}
\]

% archive sheet number: 74
\scanpage{76}
On voit directement que $\unsure{f'|\wp'}(z)$ est elliptique en y changeant $z$ en $z+\unsure{\text{deux}|2\omega_i|2\omega_1}$ \deleted{\unsure{périodes}}. De $f'(z)-f'(z+2\omega_i)=0$ on déduit $f(z)-f(z+2\omega_i)=\mathrm{c}^{\mathrm{te}}$, \unsure{dite|cette} constante est en particulier égale à $f(-\omega_i)-f(\omega_i)$, or $f$ est paire, donc la constante est nulle. \unsure{Par suite}

Il existe une fonction elliptique et une seule (de périodes $2\omega_1$ et $2\omega_2$) qui n'ait comme pôles que les $z\in\unsure{\mathfrak P|\mathfrak B}$, \deleted{\unsure{et ait pour} pôle double et partie dominante $\frac1{z^2}$.} \inserted{et telle que $\unsure{\wp|p}(z)-\frac1{z^2}$ soit holomorphe \inserted{et nulle} à l'origine.} On la note $\wp(u)$, et elle est donnée par l'expression de Mittag-Leffler \unsure{(1)|(2)}. Cette fonction est paire, \deleted{et $\wp(z)-\frac1{z^2}$ est nulle \unsure{à l'origine}} :
\[
\wp z=\frac1{z^2}+c_1z^2+c_2z^4+\ldots
\]
(L'unicité résulte immédiatement de a) ; le reste de la proposition est immédiat)

c) \uline{Toute fonction elliptique \inserted{paire} \inserted{$f$} n'ayant de pôles qu'aux points de $\mathfrak P$, est un polynôme en $\wp$ \deleted{\unsure{,}} et réciproquement. -- Ce polynôme est unique}

La réciproque est évidente, ainsi que l'unicité, $\wp$ étant transcendant sur $C$. -- $f$ aura un pôle d'ordre pair $2k$ au point $0$, nous procédons par récurrence sur $k$ (évident si $k=0$). Or, si \deleted{$f(z)\ \frac{C_k}{\unsure{z^{2k}|\wp(z)}}$} si $f(u)=\frac{C_k}{u^{2k}}+\ldots$, $f(u)-C_k\wp(u)^k$ sera \deleted{\unsure{sera} fonction elliptique ayant un pôle \unsure{plus}} une fonction elliptique paire, n'ayant \deleted{d'ordre $2(k-1)$ au point} $0$, de pôles que dans $\mathfrak P$, et $0$ pour pôle d'ordre au plus $2(k-1)$, ce qui achève la démonstration.

\uline{\unsure{Noter|Notons} que le degré de $P$, \unsure{si|où} $f=P(\wp)$, est $k$}

d) \uline{Toute fonction elliptique paire $f$ est une fonction rationnelle en $\wp$, et réciproquement}. -- \uline{unicité}

Soit $P$ le parallélogramme sous-tendu par \unsure{$f$|$\mathfrak P$}, soient $a_1\ldots a_k$ \deleted{et \unsure{répétés un nombre de}} \inserted{l'ensemble des} pôles de $f$ \inserted{dans $P$}, non dans $\mathfrak P$, \deleted{et on suppose}

% archive sheet number: 75
\scanpage{77}
\deleted{de la forme ${\unsure{\wp|p}}'R_2(\wp)$, $R_2$ fonction rationnelle en $\wp$, montrons que le dénominateur $Q$ de $R_2$ est une constante. Autrement, $Q(\wp)$\inserted{$=\prod(\wp\unsure{-a_i|-a})$} aurait des zéros hors de $\unsure{\mathfrak P|\mathfrak B}$, il faudrait que ces zéros fussent zéros de $\wp'$, avec une multiplicité au plus égale à celle dans $\wp'$ ; or $\wp'$ a}

\deleted{Prouvons d'abord le lemme}

\deleted{\unsure{h)|b)|la} \uline{Une fonction elliptique a au moins deux pôles non congrus mod $\mathfrak P$, ou un pôle double}}

\deleted{\deleted{\unsure{On peut le déduire du fait que} \unsure{toute}} application du tore dans la sphère a un degré topologique pair. Plus élémentairement, s'il n'y avait que deux pôles non congrus mod $\mathfrak P$, on pourrait par translation ramener \unsure{tous|donc} les pôles dans $\mathfrak P$. Si $0$ était pôle simple, \inserted{$f=\frac1z+\ldots$} $f^3+a\wp'$ serait fonction elliptique holomorphe, donc on aurait $f^3=\mathrm{c}^{\mathrm{te}}\wp'$ (b), \deleted{\illegible{1}}}

\deleted{\deleted{une \unsure{place} ; $f$ \unsure{paire} une fonction propre. $\unsure{f^2=\wp|f^2=P}$}}

% slanted
\margin{\deleted{On se sert du théorème : une fonction elliptique prend \unsure{le|un} $=$ nombre de fois sur $C/\mathfrak P$ n'importe quelle valeur complexe (finie ou infinie)}}

\deleted{$\wp'$ aura autant de zéros (dans $C/\mathfrak P$) \inserted{que de} pôles, savoir $3$, \deleted{et on vérifie aisément, en} d'autre part, comme \deleted{\unsure{d'autre}} \deleted{(b) que $\omega_1$ et $\omega_2$ sont ses pôles ; en \unsure{se donnant}} fonction elliptique \uline{impaire} n'ayant pas de pôles en $\omega_1,\omega_2,-\omega_3=\omega_1+\omega_2)$ $\wp'$ s'annule en les points ($f(\omega_i)=-f(-\omega_i)$, or $f(-\omega_i)=f(\omega_i)$ puisque $2\omega_i$ est une période). Il suit, \unsure{ainsi|ici} que \uline{$\wp'(u)$ \deleted{\unsure{en}} \inserted{a} $\omega_1,\omega_2,\omega_3$ comme pôles simples, (et pas d'autres dans $P$)}. $f'$ n'est donc pas \deleted{\unsure{un cube}} de la forme $f^3$.}

% archive sheet number: 76
\scanpage{78}
\deleted{mod $2\omega_1,2\omega_2$ deux à deux}. $l_i$ la multiplicité de $a_i$. Alors $g(z)=f(z)\prod_{1\leq i\leq k}\bigl[\unsure{\wp|p}(z)-\wp(a_i)\bigr]^{l_i}$ est fonction elliptique paire, n'ayant plus de pôles que dans $\unsure{\mathfrak P|\mathfrak B}$. Il suffit donc d'appliquer c).

\deleted{e) Toute fonction elliptique impaire $f$ est égale à $\wp'g$, où $g$ est une fonction elliptique paire}

e) \uline{Le corps $E$ des fonctions elliptiques est de degré 2 sur $C\langle\wp\rangle$}. En effet, $C\langle\wp\rangle$ est le corps des invariants de l'automorphisme $f\to\unsure{\check f|\breve f}$ de $E$, qui est d'ordre 2 (noter que $E\ne C\langle\wp\rangle$, car $\wp'$ est \uline{impaire})

comme corollaire, puisque $\wp'\in\complement C\langle\wp\rangle$,

f) \uline{Toute fonction elliptique $f$ se met d'une façon et d'une seule sous la forme}
\[
f=R_1(\wp)+\wp'R_2(\wp)\tag*{(3)}
\]
\uline{où $R_1$ et $R_2$ sont des fractions rationnelles}

\deleted{\unsure{Énonçons|d'ailleurs}} \inserted{Notons encore}

g) \deleted{\unsure{La totalité des pôles|Le nombre de pôles}} \uline{Pour que la fonction $f$ dans (3) ait ses pôles contenus dans $\mathfrak P$, il faut et il suffit que $R_1$ et $R_2$ soient des polynômes ; alors $0$ est au moins pôle double}.

\deleted{Il suffit de prouver que si les pôles de $f$ sont tous dans $\mathfrak P$, $f$ est de la forme $P_1(\wp)+\wp'\unsure{R_2|P_2}(\wp)$, $P_1$ et $P_2$ polynômes. Ecrivant $f=\frac{f+\check f}{2}+\frac{f-\check f}{2}$, on est ramené au cas où $f$ est paire -- traité dans c) -- resp\textsuperscript{t} où $f$ est impaire. Mais \deleted{sur} (3) on voit que cela signifie que $f$ est}

% archive sheet number: 77
\scanpage{79}
La preuve de g) est \deleted{\unsure{évidemment|immédiatement}} facile (récurrence sur la \deleted{\unsure{nombre des}}\inserted{multiplicité} du pôle\deleted{\unsure{s}} à l'origine) \unsure{si on le sait|si on la sait|puisque la suite}, est au moins d'ordre 2).

On montre maintenant \inserted{aussi} facilement que \deleted{\unsure{toute}} fonction elliptique impaire est de la forme $\mathord{\unsure{\wp|p}}'R(\wp)$ -- $R$ fraction rationnelle -- (c'est évident d'après b) etc où $R$ est un polynôme si et seulement si l'ens des pôles est $\subset\unsure{\mathfrak P|\mathfrak B}$ (appliquer g))

\noindent \unsure{2)|II)|D)} \fbox{\uline{Compléments sur $\wp,\wp'$}}

\[
\left\{\begin{aligned}
\wp&=\frac1{z^2}+\sum_{m,n}'\left(\frac1{(z-2m\omega_1-2n\omega_2)^2}-\frac1{(2m\omega_1+2n\omega_2)^2}\right)\\
\wp'&=-2\sum_{m,n}\frac1{(z-2m\omega_1-2n\omega_2)^3}\deleted{\unsure{=-\frac2{z^3}+}}
\end{aligned}\right.\tag*{(1)}
\]
\[
\wp z=\frac1{z^2}\underset{\unsure{\times}}{+}c_1z^{\unsure{2|?}}+c_2z^4+\ldots\tag*{(2)}
\]
Les $c_i$ s'obtiennent \unsure{à|aisément à} partir de (1), en dérivant les deux expressions données de $\wp z$, et en comparant pour $z\to0$
\[
\left\{\begin{aligned}
c_1&=3\sum_{m,n}'\frac1{(2m\omega_1+2n\omega_2)^4}\\
&\quad -\quad -\quad -\\
c_k&=(2k+1)\sum_{m,n}'\frac1{(2m\omega_1+2n\omega_2)^{2\unsure{n|k}+2}}
\end{aligned}\right.\tag*{(3)}
\]
Les zéros de $\wp'(z)$ dans $P$ sont $\omega_1,\omega_2,-\omega_3=\omega_1+\omega_2$. On pose \fbox{$\wp\omega_i=e_i$} $\omega_i$ est racine double de l'équation $\wp z\deleted{\unsure{+}}-e_i=0$ ; il suit que \uline{les $e_i$ sont distincts}, \unsure{car|sinon} $\wp z-e_i$ n'a que $2$ racines (distinctes ou non à priori) dans $C/\mathfrak P$

% archive sheet number: 78
\scanpage{80}
\uline{\unsure{Équation|Équations} différentielle de $\unsure{\wp|p}z$}. De façon générale, si $f$ est fonction elliptique, $f'$ est fonction elliptique, donc $f$ et $f'$ sont liés par une relation polynômiale -- \uline{toute fonction elliptique satisfait à une équation différentielle polynômiale du premier ordre}. -- Pour $\wp$, on sait que $\wp'$ est de degré $2$ sur $C\langle\wp\rangle$ ; cela peut se retrouver directement : $\wp^{\prime 2}$ est fonction elliptique paire, de degré $6$, donc polynôme en $\wp$ de degré $3$ ; on voit immédiatement \deleted{\unsure{que}} \inserted{par identification} que ce polynôme est \deleted{\unsure{de la forme}}
\[
\wp'=4\wp^3\deleted{\unsure{+}}\unsure{-}20c_1\wp-28c_2\deleted{\unsure{\wp}}
\]
Soit
\[
\boxed{\wp^{\prime 2}=4\wp^3-g_2\wp-g_3\deleted{\unsure{\wp}}}\tag*{(\unsure{7|4})}
\]
\[
(g_2=20c_1,\quad g_3=28c_2)
\]
\deleted{Le calcul (en particulier du coefficient de $\wp'$) \unsure{peut être évité}) \unsure{la relation},} en remarquant que $\frac{\wp^{\prime 2}}{(\wp-e_1)(\wp-e_2)(\wp-e_3)}$ est fonction elliptique \uline{holomorphe}, (car on vérifie qu'il n'y a pas \deleted{\unsure{de}} pôles ni à l'origine, ni aux $\unsure{\omega_i|e_i}$) donc une constante, dont on vérifie qu'elle est égale à $4$ (regarder les développements de $\wp$ et $\wp'$ à l'origine) d'où
\[
\boxed{\wp^{\prime \unsure{2|?}}=4(\wp-e_1)(\wp-e_2)(\wp-e_3)}\tag*{(5)}
\]
ce qui prouve que le polynôme en $\wp$ \unsure{cherché} est $4(\wp^3-e_1)(\wp-e_2)(\wp-e_3)$. \deleted{$\sum e_i=0$, \unsure{la différence de $\wp'$ et du second membre.}}

% archive sheet number: 79
\scanpage{81}
En poussant plus loin le calcul d'identification, on obtiendrait $c_3$ -- $c_n$ -- en fonction de $c_1,c_2$. Ce calcul est d'ailleurs simplifié considérablement en utilisant les formules suivantes, déduites de (4) par dérivation
\[
\boxed{\mathord{\unsure{\wp|p}}''=6\wp^2-\frac12g_2\qquad\wp'''=12\wp\wp'}\tag*{(6)}
\]
On obtient p. ex
\[
\boxed{c_3=\frac{c_1^2}3,\quad c_4=\frac3{11}c_1c_2}\tag*{(7)}
\]

\noindent\deleted{\unsure{2)}} 3) \fbox{\uline{Fonction $\zeta(u)$}}

\[
\zeta u-\frac1u=\int_0^u-\left(\wp z-\frac1{z^2}\right)dz\qquad\zeta'u=-\wp u\tag*{(7)}
\]
\[
\zeta u=\frac1u-\frac{c_2}{3}u^3\ldots-\frac{c_n}{2n+1}u^{2n+1}\ldots-\tag*{(8)}
\]
$\zeta$ est impaire, admet les périodes pour pôles simples, avec \deleted{\unsure{pôles}} \unsure{résidu|résidus} 1.
\[
\zeta z=\frac1z+\sum_{\omega\in{\unsure{\mathfrak P|\mathfrak B}}^{*}}\left(\frac1{z\unsure{-|+}\omega}+\frac1{\omega}+\frac{z}{\omega^2}\right)\tag*{(9)}
\]
$\zeta$ n'est pas elliptique (\unsure{voir plus loin}). --
\[
\boxed{\zeta(z+2\omega)-\zeta(z)=2\eta}\tag*{\textcircled{10}}
\]
($=\mathrm c^{\mathrm{te}}$ ; comme on voit en dérivant le premier membre).

Faisant $z=-\omega$, on obtient
\[
\eta=\zeta(\omega)\tag*{\textcircled{11}}
\]
; $\eta_1,\eta_2,\omega_1,\omega_2$ sont liés par une relation ; si $\omega_1,\omega_2$ sont en disposition directe, on a
\[
\eta_1\omega_2-\eta_2\omega_1=\frac{i\pi}{\unsure{2}}
\]

% archive sheet number: 80
\scanpage{82}
(en calculant $\int z\unsure{\wp|p}z$ le long d'un parallélogramme de périodes ayant pour centre l'origine).

\uline{Formule d'Hermite}. Si la fonction elliptique $f$ a, dans le parallélogramme des périodes $P$, un syst fondamental de pôles $a_i$, de parties principales
\[
\frac{A_{i,1}}{z-a_i}+\frac{A_{i,2}}{(z-a_i)^2}\ \unsure{\ldots|{-}}\ +\frac{A_{i,k_i}}{(z-a_i)^{k_i}}
\]
on a
\[
\deleted{f(z)=\mathrm c^{\mathrm{te}}+\sum_i\bigl(A_{i,1}\zeta(z-a_i)+A_{i,2}\wp(z-a_i)\bigr)}
\]
\[
\begin{aligned}
f(z)&=\mathrm c^{\mathrm{te}}+\sum_i\biggl(A_{i,1}\zeta(z-a_i)+A_{i,2}\wp(z-a_i)+\ldots\\
&\qquad+\frac{(-1)^{\unsure{k_i|k}}}{\unsure{(k_i-1)!|(k_i-2)!}}A_{i,k_i}\wp^{(\unsure{k_i|k})}(z-a_i)\biggr)
\end{aligned}\tag*{\textcircled{12}}
\]
(le deuxième membre est elliptique, grâce à $\sum A_i=0$)

\noindent 4) \fbox{Fonction $\sigma z$}
\[
\sigma z=z\prod_{\omega\in{\unsure{\mathfrak P|\mathfrak B}}^{*}}\left(1-\frac z\omega\right)e^{\frac z\omega+\frac{z^2}{2\omega^2}}\tag*{\textcircled{13}}
\]
\[
\frac{\sigma'z}{\sigma z}=(\lg\sigma z)'=\zeta z\qquad \sigma z=ze^{\int_0^z\left(\zeta u-\frac1u\right)du}\tag*{\textcircled{14}}
\]
$\sigma$ fonction entière impaire, ayant les $\omega\in\mathfrak P$ pour \deleted{\unsure{pôles}} zéros simples (et pas d'autres zéros)
\[
\sigma u=u-c_1\frac{u^5}{12}-c_2\frac{u^7}{30}\tag*{\textcircled{15}}
\]

% archive sheet number: 81
